% No subheadings

\section*{Introduction} % 1000 - 1500 words
%% Intro %% 
% In recent years, the rapid growth of artificial intelligence (AI) and machine learning (ML) has resulted in increasingly complex and large models, leading to substantial latency and energy consumption challenges that can make edge deployment too slow and power-hungry, and cloud deployment commercially impractical. These issues are further magnified by the decelerating performance gains of conventional computing architectures as they approach fundamental physical limits, the imperative for energy-efficient edge-computing devices to accommodate the swift expansion of the Internet of Things (IoT), and the necessity for real-time systems capable of functioning in closed-loop environments. As a result, the urgency for alternative computational paradigms has intensified. 



In recent years, artificial intelligence (AI) and machine learning (ML) models have grown increasingly large and complex in pursuit of higher accuracy and capabilities~\cite{?}. However, larger models lead to substantial latency and energy consumption challenges that can render edge deployment infeasible due to slow speed and high power requirements. Moreover, cloud deployment of massive models is often commercially impractical due to high compute costs~\cite{?}. These issues are compounded by the slowing performance gains in conventional computing architectures as they reach physical limits of efficiency and speed~\cite{?}. At the same time, the expanding Internet of Things (IoT) necessitates low-power edge devices with AI capabilities~\cite{?}. 

Neuromorphic computing has emerged as a promising area in addressing these challenges, aiming to unlock key hallmarks of biological intelligence by porting primitives and computational strategies employed in the brain into engineered computing devices and methods~\cite{Schuman2022, james17historical, thakur2018largescale}.
Neuromorphic systems hold a critical position in the investigation of novel architectures, as the brain exemplifies an exceptional model for accomplishing scalable, energy-efficient, and real-time embodied computation.

%% Background %%
Originally, the term ``neuromorphic'' specifically referred to approaches aiming to mimic the biophysics of the brain by leveraging the physical properties of silicon, as first proposed by Mead in the 1980s~\cite{Mead:1990hw}. However, neuromorphic computing research has since expanded to encompass a wide variety of brain-inspired techniques at the algorithmic, hardware, and system levels~\cite{schuman2017survey}. While diverse, neuromorphic approaches generally utilize mechanisms that emulate biological processes more closely than conventional methods. This bio-realistic modeling aims to reproduce high-level performance characteristics of biological neural systems. At the algorithmic level, neuromorphic techniques seek to achieve expanded learning capabilities, predictive intelligence, data efficiency, and adaptive learning akin to the brain through the use of spiking neurons, plastic synapses, and neural dynamics. In neuromorphic hardware, the goal is to capture the energy efficiency and real-time processing of neural systems by implementing architectures based on primitives such as sparse, event-driven, and distributed computation.

% Initially, the term \textit{neuromorphic} referred specifically to approaches that aimed to emulate the biophysics of the brain through the use of the physical properties of silicon substrate, as proposed by Mead in the 1980's~\cite{Mead:1990hw}. However, the field of neuromorphic computing research has since grown to encompass a wide range of brain-inspired computing techniques~\cite{schuman2017survey} at the algorithmic, hardware, and co-designed system levels. 
% While the range of approaches is diverse, neuromorphic computing research generally utilizes mechanisms emulating or simulating biophysical properties to a greater degree than conventional methods, in order to demonstrate shared high-level performance goals. Algorithmic approaches strive towards using such mechanisms with goals of expanded learning and processing capabilities, such as higher predictive intelligence, improved data efficiency, and adaptive learning. Systems of neuromorphic algorithms deployed to hardware seek greater energy efficiency and temporal awareness than conventional hardware.

Neuromorphic algorithms encompass emerging techniques inspired by neuroscience that can ultimately be accelerated by leveraging the properties of neuromorphic hardware. A prime example is spiking neural networks (SNNs), which incorporate biologically-realistic spiking neurons and synaptic plasticity~\cite{?}. In the early stages of hardware development, neuromorphic algorithms may not be fully supported by dedicated neuromorphic chips due to limitations in scale or attributes. As a result, current research relies heavily on simulating these algorithms on conventional hardware like CPUs and GPUs, with the goal of eventual deployment on specialized neuromorphic platforms once capabilities advance.

In recent years, the advent of neuromorphic chips has led to brain-inspired computing in hardware through diverse approaches, including analog circuits mimicking neural dynamics, digital spiking networks, and non von-Neumann architectures~\cite{?}. Key properties of these chips include event-driven spike-based processing~\cite{?}, in-memory computing with emerging devices~\cite{?}, sparse and low-precision computation~\cite{?}, as well as adaptive and fault-tolerant characteristics~\cite{?}. These features are intended to reproduce neural systems' efficiency, real-time responsiveness, and resilience. Neuromorphic chips target a wide range of applications, from neuroscience research to low-power intelligence at the edge to datacenter acceleration of AI workloads.



% Neuromorphic \textit{algorithms} encompass emerging neuroscience-inspired algorithms and methods that can be ultimately accelerated utilizing the sparse, event-driven, temporal, and distributed nature of neuromorphic hardware, such as spiking neural networks (SNNs). Depending on their development stage, the attributes or scale of neuromorphic algorithms may not yet be adequately supported by existing neuromorphic hardware. As such, current algorithm exploration often makes use of simulated execution on readily-available conventional hardware such as central processing units (CPUs) and graphics processing units (GPUs), with the eventual goal of deployment to neuromorphic systems. 

% Neuromorphic \textit{systems} seek to realize biologically-inspired computing machines through a variety of hardware approaches, such as analog emulation and digital simulation, spike- or event-based computation and communication, non-von-Neumann architectures, near- and in-memory processing with emerging memory devices, as well as various properties such as low-resolution, sparse, noisy, and adaptive processing. Such machines are intended for a wide range of use-cases, from neuroscientific bio-emulation research, to low-power edge intelligence, to datacenter-scale acceleration~\cite{frenkel23bottom}.

While possessing promise, it remains unclear which neuromorphic computing approaches truly deliver on anticipated capabilities due to insufficient benchmarking practices for fair assessment and comparison. Without widely adopted standard benchmarks, progress is impeded by a lack of clarity on which solutions actually achieve breakthrough efficiency, speed, and intelligence ~\cite{Davies19_bencprog, Schuman2022}. Although some application-specific benchmarks have emerged, such as for vision, audition, and simulation tasks, these narrow efforts do not adequately evaluate real-world viability across problem domains~\cite{orchard2015converting,Amir_etal17_lowpowe,Cramer_2022,Ostrau2022,closedloop-andre}. 

Furthermore, existing benchmarks are rarely implemented in a consistent manner in the scientific literature, limiting replicability. Current benchmarking gaps severely hinder the ability to discern which proposed techniques live up to their purported potential. Uncertainty around true capabilities risks wasting effort and resources on approaches that may not deliver fundamentally new computing paradigms in practice. Only through rigorous and standardized benchmarking can the neuromorphic community objectively assess promised versus actual achievements, and make evidence-based decisions on which directions show genuine promise versus hype. Establishing widely adopted benchmark suites to evaluate real-world viability will enable clearer discrimination of successful versus unsuccessful approaches on key criteria like efficiency, speed, accuracy, and scalability. This clarity is essential to focus research and commercialization efforts on techniques that concretely demonstrate enhanced functionality and performance relative to conventional computing.
However, developing and establishing such standardized neuromorphic benchmarks requires us to address the following challenges:

%% Challenges/Problem statement %%
% Despite its promises, progress in the field of neuromorphic research is impeded by a lack of fair and widely-used objective metrics and benchmarks, which has been documented in numerous calls to action~\cite{Davies19_bencprog, Schuman2022}.
% Previously, other neuromorphic benchmarks have also been proposed, from classical vision~\cite{orchard2015converting,Amir_etal17_lowpowe} and audition tasks~\cite{Cramer_2022}, to open-loop~\cite{Ostrau2022} and closed-loop~\cite{closedloop-andre} tasks, or on the performance of SNN simulators~\cite{KULKARNI2021}.
% However, these works do not address key existing challenges in adopting neuromorphic benchmarks:

\begin{itemize}
    \item \textbf{Lack of a formal definition.} The lack of a formal definition for what constitutes a "neuromorphic" solution makes it difficult to delineate precise criteria for benchmarking. Given the range of techniques from simulations to custom neural hardware, benchmarks must remain inclusive to different abstraction levels and goals. Overly narrow benchmarks risk unfairly excluding promising methods or imposing assumptions. Ideal benchmarks should focus on capabilities of interest like temporal processing and efficiency while allowing flexible implementation. Comparisons to conventional techniques are also key for demonstrating advantages.

    \item \textbf{Implementation diversity.} Implementation diversity has been instrumental for exploring bio-inspired paradigms but creates portability issues. The myriad of frameworks and tools follow different methodologies based on intended use cases, from modeling neuroscience principles~\cite{nest2007} to automating neural network design~\cite{eshraghian2021training}. However, variation in independent designs hampers standardizing benchmark methodology. This limits adoption and actionable insights. Portable and unified benchmarks are needed to fairly compare disparate approaches.

    
    \item \textbf{Rapid research evolution.} Rapid evolution means benchmarks must iterate alongside advancing research. As breakthrough techniques continually emerge, evaluation practices must expand to track progress over time. Structured benchmark versions can incorporate new developments while maintaining continuity and backward compatibility. This iterativeness enables benchmarking to foster innovation rather than impede it.
    
\end{itemize}

% \begin{itemize}
%     \item \textbf{Lack of a formal definition.} The rich variety of approaches to exploring brain-inspired principles creates difficulties in defining a closed set of criteria for what should be benchmarked as a \textit{neuromorphic} solution.
%     This challenge necessitates \textbf{\textit{inclusive}} benchmarks that can be applied generally across the spectrum of potential approaches, while being composed of tasks and metrics which are of interest to neuromorphic research (e.g., temporal applications and efficiency metrics). This way, the benchmark itself does not present obstacles to neuromorphic assessment by making assumptions of solutions.    
%     Furthermore, inclusive benchmarks that can apply to non-neuromorphic solutions as well are necessary for demonstrating the advantages of neuromorphic approaches in direct comparison with conventional approaches.
    
%     \item \textbf{Implementation diversity.} A wide array of different frameworks, which in many cases aim towards different goals (e.g., neuroscientific simulation~\cite{nest2007} and automatic SNN configuration~\cite{eshraghian2021training}), are used in neuromorphic research. This diversity has been instrumental in exploring the landscape of bio-inspired techniques following different methodologies and abstraction levels, but the broad variation of framework goals and independent implementation styles create barriers in portability and standardization of benchmark methodology, which in turn limits the \textbf{\textit{actionable}} ease of benchmark implementation.
    
%     \item \textbf{Rapid research evolution.} As part of an emerging field, neuromorphic approaches continually and rapidly evolve. As the research community continues to make technological progress, so too should benchmark suites and methodology expand to foster inclusion and capture salient performance metrics. An \textbf{\textit{iterative}} benchmark framework with structured successive versions will facilitate productive foundational and incrementally improving performance evaluation.
% \end{itemize}

To address current challenges in evaluating neuromorphic solutions, this article presents NeuroBench - a new multi-task benchmark framework developed collaboratively across academia and industry. NeuroBench aims to advance neuromorphic benchmarking in three key ways. First, it utilizes general task definitions and hierarchical performance metrics relevant to diverse techniques. This inclusive design minimizes assumptions and enables participation beyond specialized neuromorphic approaches. Second, NeuroBench provides an open-source harness tool to standardize benchmark implementation across frameworks, facilitating actionable comparisons. Finally, establishing NeuroBench as a continuous community initiative enables iterative evolution akin to MLPerf~\cite{Reddi2020,mattson2020mlperf}, ensuring adaptation alongside advancing research.

NeuroBench consists of two complementary tracks - an algorithm track for hardware-independent evaluation and a system track for fully deployed solutions. 
The algorithm track defines four novel benchmarks for neuromorphic methods across diverse domains of few-shot continual learning, computer vision, motor cortical emulation, and chaotic forecasting, and utilizes complexity metrics to analyze solution costs. Such hardware-independent benchmarking enables algorithmic exploration and prototyping, especially when simulating execution on non-neuromorphic platforms. Meanwhile, the system track defines standard protocols to measure the real-world speed and efficiency of neuromorphic hardware on benchmarks ranging from standard machine learning tasks to promising fields for neuromorphic systems such as optimization.
Both the algorithm and system track will be extended and co-developed as NeuroBench continues to expand.
\JY{How to note that the system track doesn't have anything completed as of this article?}
% This separation supports algorithm prototyping while also benchmarking hardware performance. Both tracks will be extended over time through collaborative refinement. In summary, NeuroBench seeks to drive neuromorphic progress through inclusive scope, standardized implementation, and community-driven iteration.

% To tackle these challenges, this article presents NeuroBench, a multi-task benchmark framework, collaboratively developed by academic and industry partners with a stake in neuromorphic computing solutions.
% NeuroBench addresses the existing benchmark challenges by advancing prior work in three distinct ways. 
% Firstly, the benchmark framework reduces assumptions regarding the specific solution being assessed, encouraging inclusive participation of neuromorphic and non-neuromorphic approaches by utilizing general, task-level benchmarking and hierarchical metric definitions which capture key performance indicators of interest to the neuromorphic community. Secondly, the NeuroBench benchmarks are associated with a common open-source benchmark harness tool which facilitates actionable benchmark implementation and offers structure for further expansion to neuromorphic algorithm frameworks and systems. Thirdly, NeuroBench establishes a continuous, community-driven initiative designed to evolve over time to ensure representation and relevance to neuromorphic research, analogous to the successful MLPerf benchmark framework~\cite{Reddi2020}. 

% NeuroBench consists of two tracks: an \textit{algorithm track} that addresses algorithmic solutions in a hardware-independent manner, and a \textit{system track} that tackles full-system solutions. The dual-track structure offers a systematic separation of benchmarking for algorithm exploration and prototyping, especially when simulating execution on non-neuromorphic platforms, and benchmarking of neuromorphic hardware performance against existing computing hardware. 
% The tracks will be extended and co-developed as NeuroBench continues to be refined.

% This article presents the framework of the NeuroBench benchmark suite. 
The following section organizes descriptions of the algorithm track benchmark framework and its baseline results, as well as specifications of the system track benchmark framework and tasks. Further details regarding the benchmark metric formulations, task specifications, and baseline solutions can be found in the Methods section. Baseline results on NeuroBench benchmarks outline unexplored research opportunities in optimizing algorithmic architectures and training of sparse, stateful models to achieve greater performance and resource efficiency.
As NeuroBench is intended to continually grow over time, the latest developments and opportunities to engage with the project are reported on the website\footnote{\href{https://neurobench.ai}{https://neurobench.ai}}.